\section{Lowering and numerics}

The lowering from the tensor level to the instruction level is a representation
change, and the dialect conversion framework is the right tool for it, but the
framework's materialization rules took care to get right. A type converter maps
every tensor to a scratchpad buffer, and target and source materializations insert
a load or a store at the boundary. The subtlety is keeping DRAM values, the
function arguments and the constants, as tensors while the intermediates become
buffers, so that DMA appears only where the two memories meet. A constant becomes
a DRAM constant followed by a load, which makes the boundary explicit.

The one genuine correctness bug appeared only under a tight scratchpad budget,
and the benchmark suite caught it. At a small budget the allocator spills a buffer
by storing it to DRAM after it is produced and reloading it before its next use.
The maximum error against onnxruntime jumped from the fp32 rounding level to about
eight hundredths. The cause was in the instruction encoder, which assigned DRAM
offsets only to inputs, constants, and returned values. A spill store's result is
a DRAM temporary that is none of those, so with no offset assigned it defaulted to
offset zero and overwrote the model input, and the reload read the input back
instead of the spilled buffer. The scratchpad allocation test had checked only
that spilling inserted the right number of DMA instructions, never the numerics,
so it did not catch this. The fix gives every spill temporary its own DRAM region,
and an end to end test at a tight budget now locks the numerics in.
